\section{Generalized Reeb orbits on polytopes}
\label{sec:generalized-reeb-orbits-polytopes}

On the interior of a polytope facet, the Reeb velocity is constant. At an
intersection of facets, the generalized characteristic equation permits
convex combinations of their velocities. Haim--Kislev's theorem nevertheless
guarantees a capacity minimizer that uses only the individual facet
velocities, each at most once. Its action is therefore determined by a finite
ordered list of facets and the time spent in each direction.

We first derive this finite description and distinguish closure of the curve
from its placement on the boundary. We then prove the existence theorem by
rearranging curves in the dual variational problem.


\subsection{Facet velocities}\label{subsec:generalized-reeb-orbits-polytopes-definition}

Let
\[
  K=\{x\in\mathbb{R}^4:\langle a_i,x\rangle\leq 1,\; i=1,\ldots,F\}
\]
be a full-dimensional polytope with \(0\in\operatorname{int}(K)\). We assume
that the halfspace description is irredundant, so the facet corresponding to
\(a_i\) is
\[
  F_i=\{x\in K:\langle a_i,x\rangle=1\}.
\]
The normalized halfspace description gives
\[
  g_K(x)=\max_{1\leq i\leq F}\langle a_i,x\rangle.
\]
If \(x\in\partial K\), let
\(I(x)=\{i:x\in F_i\}\) be its active facets.  The subdifferential rules for a
finite maximum and for the square on the level set \(g_K=1\) give
\begin{equation}\label{eq:generalized-reeb-polytope-subdifferential}
  \partial H_K(x)
  =2\operatorname{conv}\{a_i:i\in I(x)\}.
\end{equation}
Thus Definition~\ref{def:preliminaries-generalized-characteristic} becomes a
finite differential inclusion.  In this normalization the pure
contact-normalized Reeb direction on \(F_i\) is
\begin{equation}\label{eq:generalized-reeb-facet-vector}
  R_i=2J_0a_i.
\end{equation}
Indeed \(R_i\) is tangent to \(F_i\), because
\(\langle a_i,J_0a_i\rangle=0\), and for \(x\in F_i\),
\[
  \lambda_0|_x(R_i)
  =\frac{1}{2}\langle -J_0R_i,x\rangle
  =\langle a_i,x\rangle
  =1.
\]

For a Lagrangian product, this formula makes the exchange between the two
factors explicit.  A row \(a_i=(u,0)\) from a facet of the \(q\)-factor and a
row \(a_j=(0,v)\) from a facet of the \(p\)-factor give
\[
  R_i=2J_0(u,0)=(0,2u),
  \qquad
  R_j=2J_0(0,v)=(-2v,0).
\]
Thus a segment on a \(q\)-facet moves only in the \(p\)-coordinates, while a
segment on a \(p\)-facet moves only in the \(q\)-coordinates.

A \emph{generalized Reeb orbit} on \(K\) is a contact-normalized generalized
characteristic in the sense of
Definition~\ref{def:preliminaries-generalized-characteristic}.  By
\eqref{eq:generalized-reeb-polytope-subdifferential}, this means precisely that
\(\gamma\in W^{1,2}([0,T],\mathbb{R}^4)\) is closed, has image in
\(\partial K\), and satisfies
\begin{equation}\label{eq:generalized-reeb-inclusion}
  \dot{\gamma}(t)
  \in \operatorname{conv}\{R_i:\gamma(t)\in F_i\}
  \quad\text{for a.e. }t\in[0,T].
\end{equation}
The general normalization already shows that such an orbit has
\[
  A(\gamma)=T.
\]

We call such an orbit \emph{simple} if, after choosing its starting point
at a velocity change, it consists of finitely many segments with velocities
$R_i$, and each $R_i$ occurs in at most one segment. Here simplicity refers
to the use of facet directions. In particular, it includes the requirement
that the velocities be individual $R_i$, rather than nontrivial convex
combinations of them.

\subsection{Words, dwell times, and closure}
\label{subsec:generalized-reeb-orbits-polytopes-words}

A simple orbit is described by the ordered facet labels it uses, the time
spent in each direction, and a starting point. The \emph{active word}
contains only directions used for positive time. A cyclic word
\[
  \sigma=(\sigma_1,\ldots,\sigma_m)
\]
with distinct entries, together with dwell times \(\tau_k>0\) and a base point
\(b=\gamma(0)\), describes the velocity data
\[
  \dot{\gamma}(t)=R_{\sigma_k}
  \quad\text{for a time interval of length }\tau_k.
\]
Such data can represent a closed orbit only if the closing equation holds:
\begin{equation}\label{eq:generalized-reeb-closure}
  \sum_{k=1}^m \tau_k R_{\sigma_k}=0.
\end{equation}
This equation closes the broken path. A further condition is needed to
place each segment on its prescribed facet: there must be a suitable
starting point $b$. Lemma~\ref{lem:generalized-reeb-base-point-recovery}
expresses that condition by linear equalities and inequalities.


For example, closure alone fails to give an orbit in the
planar square \([-1,1]^2\), the right and left facet rows give opposite
directions \(R_r=(0,2)\) and \(R_\ell=(0,-2)\).  Equal dwell times satisfy
\(\tau R_r+\tau R_\ell=0\), but the breakpoint between the two segments would
have to lie in the disjoint facets \(F_r\) and \(F_\ell\).  No translation can
realize this closed two-segment path on the prescribed facets.

The period and action of a closed pure-word orbit are
\[
  T=\sum_{k=1}^m \tau_k,
  \qquad
  A(\gamma)=T.
\]


Cyclically shifting the facet--duration pairs and moving the starting point
to the corresponding breakpoint describes the same orbit. We therefore regard
the word as cyclic. A zero duration, when allowed in a limiting argument,
is simply deleted from the active word.

The action of a finite word can be calculated directly from its velocities and
dwell times.

\begin{lemma}[Symplectic shoelace formula]\label{lem:preliminaries-shoelace}
Let \(z\colon[0,T]\to\mathbb{R}^4\) be a closed piecewise affine curve. Suppose
that \(z\) has constant velocity \(w_k\) on consecutive intervals of length
\(\ell_k\), for \(k=1,\ldots,m\). Then
\begin{equation}\label{eq:preliminaries-shoelace}
  2A(z)
  =
  \sum_{1\leq j<k\leq m}
    \ell_j\ell_k\,\omega_0(w_j,w_k).
\end{equation}
In particular, the action of a closed pure-word Reeb orbit depends only on the
cyclic velocity word and the dwell times, not on the base point.
\end{lemma}

\begin{proof}
Write \(t_k=\sum_{j<k}\ell_j\). On the \(k\)-th interval,
\[
  z(t)=z(0)+\sum_{j<k}\ell_jw_j+(t-t_k)w_k.
\]
The contribution of \(z(0)\) to
\(\int_0^T\langle -J_0\dot z,z\rangle\,dt\) vanishes because \(z\) is closed,
so \(\sum_k\ell_kw_k=0\). The self-term
\(\langle -J_0w_k,w_k\rangle\) vanishes by skew-symmetry. Hence
\[
  2A(z)
  =
  \sum_k
    \ell_k\left\langle -J_0w_k,\sum_{j<k}\ell_jw_j\right\rangle
  =
  \sum_{j<k}\ell_j\ell_k\,\omega_0(w_j,w_k),
\]
using \(\omega_0(u,v)=\langle J_0u,v\rangle\).
\end{proof}

Breakpoints of \(\dot{\gamma}\) are velocity changes, not necessarily changes
in the minimal face containing \(\gamma(t)\): an orbit may move inside a
nonsmooth face while switching from one active facet direction to another.

\begin{lemma}[Base-point recovery]\label{lem:generalized-reeb-base-point-recovery}
Fix an active word \(\sigma=(\sigma_1,\ldots,\sigma_m)\) and dwell times
\(\tau_k>0\) satisfying the closure equation
\eqref{eq:generalized-reeb-closure}. Put
\[
  v_1=0,
  \qquad
  v_k=\sum_{j=1}^{k-1}\tau_jR_{\sigma_j}
  \quad (k=2,\ldots,m+1).
\]
Then \(v_{m+1}=0\). A point \(b\in\mathbb{R}^4\) recovers a simple generalized Reeb orbit
with these velocity data, by
\[
  \gamma(t)=b+v_k+sR_{\sigma_k},
  \qquad 0\leq s\leq \tau_k
\]
on the \(k\)-th segment, if and only if
\begin{equation}\label{eq:generalized-reeb-base-point-equalities}
  \langle a_{\sigma_k}, b+v_k\rangle=1
  \quad\text{for }k=1,\ldots,m,
\end{equation}
and
\begin{equation}\label{eq:generalized-reeb-base-point-inequalities}
  \langle a_i,b+v_k\rangle\leq 1
  \quad\text{for }i=1,\ldots,F\text{ and }k=1,\ldots,m.
\end{equation}
Equivalently, the equalities have the form
\[
  N_\sigma b = r_\sigma,
  \qquad
  (N_\sigma)_{k\bullet}=a_{\sigma_k}^{\mathsf T},
  \qquad
  (r_\sigma)_k=1-\langle a_{\sigma_k},v_k\rangle .
\]
When this system is nonempty, its solution space is an affine space of
dimension \(4-\operatorname{rank}N_\sigma\). The inequalities cut out the
valid base points as a convex polytope inside this affine space. In particular,
the set of valid base points is bounded, because \(v_1=0\) and
\eqref{eq:generalized-reeb-base-point-inequalities} imply \(b\in K\).
\end{lemma}

\begin{proof}
On the \(k\)-th segment, the prescribed facet equation is
\[
  \langle a_{\sigma_k}, b+v_k+sR_{\sigma_k}\rangle=1.
\]
The \(s\)-term vanishes because
\(\langle a_{\sigma_k},R_{\sigma_k}\rangle
=2\langle a_{\sigma_k},J_0a_{\sigma_k}\rangle=0\). Hence the whole segment
lies on \(F_{\sigma_k}\) exactly when
\eqref{eq:generalized-reeb-base-point-equalities} holds.

The segment lies in \(K\) exactly when
\[
  \langle a_i,b+v_k+sR_{\sigma_k}\rangle\leq 1
  \quad\text{for all }i\text{ and all }s\in[0,\tau_k].
\]
For fixed \(i,k\) this expression is affine in \(s\), so it is enough to check
the endpoints. The endpoint at \(s=0\) is \(b+v_k\), and the endpoint at
\(s=\tau_k\) is \(b+v_{k+1}\). As \(k\) runs cyclically and \(v_{m+1}=v_1=0\),
these endpoint checks are exactly
\eqref{eq:generalized-reeb-base-point-inequalities}. With these equalities and
inequalities, the resulting closed broken path lies on \(\partial K\) and has
velocity \(R_{\sigma_k}\) on the \(k\)-th segment, hence is a simple Reeb orbit
with the prescribed data.
\end{proof}

\begin{remark}[Non-uniqueness and computation]
The base point is not part of the finite action data. If
\(\operatorname{rank}N_\sigma=4\), the facet equations determine it uniquely
when they are consistent. If \(\operatorname{rank}N_\sigma<4\), the inequalities
may leave a positive-dimensional set of base points. This non-uniqueness does
not affect the action, which is determined by the closed velocity polygon via
Lemma~\ref{lem:preliminaries-shoelace}. Computationally, base-point recovery is
therefore a finite linear feasibility check after the word and dwell times have
been chosen.
\end{remark}

\subsection{Existence of simple minimizers}
\label{subsec:generalized-reeb-orbits-polytopes-existence-of-simple-minimizers}

The next theorem is the reason a finite search can compute capacity.
It is Haim--Kislev's simple-minimizer theorem~\cite[Theorem~1.5]{HK2017},
written with the contact-normalized velocities $R_i=2J_0a_i$.

\begin{theorem}[Simple minimizer]\label{thm:generalized-reeb-simple-minimizer}
Let \(K\subset\mathbb{R}^4\) be a full-dimensional convex polytope with
\(0\in\operatorname{int}(K)\). Among all generalized Reeb orbits on
\(\partial K\), there exists a minimum-action orbit that is simple in the
contact-normalized convention above. Hence
\[
  c_{\mathrm{EHZ}}(K)
  =
  \min\{A(\gamma):\gamma
  \text{ is a simple Reeb orbit on }\partial K\}.
\]
\end{theorem}

The difficulty is that splitting or rearranging velocities need not leave
a curve on $\partial K$. We make those changes to closed curves in the dual
problem instead. The changes preserve the total time and can be chosen to
increase action. Rescaling all durations restores the action--period
constraint. Dual minimality controls the possible action change, and a
convergent subsequence then supplies a minimizer with the desired finite form.

During the proof, a \emph{simple pure-velocity dual loop} means a closed
piecewise affine curve using each $R_i$ on at most one cyclic interval.
It need not lie on the boundary. The next five lemmas give the approximation,
rearrangement and compactness steps.

\begin{lemma}[Piecewise affine approximation]\label{lem:generalized-reeb-pw-affine-approx}
Let \(z\in W^{1,2}([0,T],\mathbb{R}^4)\) be closed and satisfy
\[
  \dot z(t)\in\operatorname{conv}\{R_1,\ldots,R_F\}
  \quad\text{for a.e. }t.
\]
For every \(\varepsilon>0\), there is a closed piecewise affine curve \(z'\)
with
\[
  \dot z'(t)\in\operatorname{conv}\{R_1,\ldots,R_F\}
\]
on each affine piece and \(\|z'-z\|_{W^{1,2}}<\varepsilon\). Consequently
\(A(z')\to A(z)\) and \(I_K(z')\to I_K(z)\) as \(\varepsilon\to0\).
\end{lemma}

\begin{proof}
Let \(C=\operatorname{conv}\{R_1,\ldots,R_F\}\). Averages over intervals of refining uniform partitions converge to
\(\dot z\) in \(L^2\). Choose a partition \(0=t_0<\cdots<t_m=T\) so that the conditional averages of
\(\dot z\) on the partition intervals are \(L^2\)-close to \(\dot z\). On
\([t_k,t_{k+1}]\), put
\[
  w_k=\frac{1}{t_{k+1}-t_k}\int_{t_k}^{t_{k+1}}\dot z(s)\,ds.
\]
Because \(C\) is convex and closed, each \(w_k\) lies in \(C\). Define \(z'\)
by \(z'(0)=z(0)\) and \(\dot z'=w_k\) on \([t_k,t_{k+1}]\). Then
\[
  z'(T)-z'(0)=\int_0^T\dot z'(t)\,dt=\int_0^T\dot z(t)\,dt=0,
\]
so \(z'\) is closed. Since \(\dot z'\to\dot z\) in \(L^2\), integration gives
\(z'\to z\) uniformly, hence in \(W^{1,2}\). The convergence of \(A\) follows
from bilinearity in \((z,\dot z)\). Both \(-J_0\dot z'\) and
\(-J_0\dot z\) take values in the compact set \(-J_0C\), on which
\(h_K^2\) is Lipschitz; hence the
\(L^2\)-convergence also gives \(I_K(z')\to I_K(z)\).
\end{proof}

\begin{lemma}[Splitting mixed velocities]\label{lem:generalized-reeb-splitting}
Let \(z\) be a closed piecewise affine curve of period \(T\) whose velocities
lie in \(\operatorname{conv}\{R_1,\ldots,R_F\}\). Then there is a closed
piecewise affine curve \(z'\), with the same period, whose velocity on each
segment is one pure Reeb vector \(R_i\), such that \(A(z')\geq A(z)\) and
\[
  I_K(z')=T.
\]
\end{lemma}

\begin{proof}
On a segment with velocity \(v=\sum_i\theta_iR_i\), \(\theta_i\ge0\) and
\(\sum_i\theta_i=1\), replace the segment by subsegments with velocities
\(R_i\) and lengths \(\theta_i\) times the original length. The total
displacement and period are unchanged. The only change in the action is the
sum of the new internal cross-terms from
Lemma~\ref{lem:preliminaries-shoelace}. Reversing the order of the inserted
subsegments changes the sign of this internal contribution, so one of the two
orders does not decrease the action.

For a pure Reeb velocity \(R_i=2J_0a_i\),
\[
  -J_0R_i=2a_i,
  \qquad
  h_K(-J_0R_i)=2h_K(a_i)=2,
\]
because \(F_i\neq\varnothing\) and \(\max_{x\in K}\langle a_i,x\rangle=1\).
Thus \(h_K^2(-J_0R_i)=4\) on every pure segment, and
\[
  I_K(z')=\frac14\int_0^T4\,dt=T.
\]
\end{proof}

\begin{lemma}[Merging repeated velocities]\label{lem:generalized-reeb-merging}
Let \(z\) be a closed piecewise affine curve whose velocity on each segment is
one of the \(R_i\). If some \(R_i\) occurs in two non-adjacent segments, then
there is a rearrangement \(z'\) that merges those two occurrences into one
contiguous segment and satisfies
\[
  A(z')\geq A(z),
  \qquad
  I_K(z')=I_K(z).
\]
Repeating this operation, and merging adjacent equal-velocity intervals,
gives a closed piecewise affine curve in which each pure Reeb vector occurs in
at most one interval. Here adjacency is understood cyclically, since the curve
is closed.
\end{lemma}

\begin{proof}
Suppose a velocity \(v\) occurs in two segments \(I_r\) and \(I_s\), with
middle segments \(M=\{r+1,\ldots,s-1\}\). Here \(I_i\) is a time interval,
\(|I_i|\) is its duration, and \(w_i\) is the velocity on that interval.
Merging \(I_s\) forward into \(I_r\)
changes only the action terms involving \(v\) and the middle velocities, and
the change is
\[
  \Delta A_{\mathrm{fwd}}
  =
  |I_s|\sum_{i\in M}|I_i|\,\omega_0(v,w_i).
\]
Merging \(I_r\) backward into \(I_s\) gives
\[
  \Delta A_{\mathrm{bwd}}
  =
  -|I_r|\sum_{i\in M}|I_i|\,\omega_0(v,w_i).
\]
The two changes have opposite signs, or both vanish, so one merge has
nonnegative action change. Closure is preserved because the total duration assigned to each
velocity is unchanged, and \(I_K\) is unchanged because it
depends only on the velocities and their durations.
\end{proof}

\begin{lemma}[Rescaling]\label{lem:generalized-reeb-rescaling}
Let \(z\) be a closed piecewise affine curve with pure Reeb velocities and
period \(T\). If all durations are multiplied by \(\beta>0\), the resulting
curve \(z'\) satisfies
\[
  A(z')=\beta^2A(z),
  \qquad
  I_K(z')=\beta I_K(z).
\]
\end{lemma}

\begin{proof}
In Lemma~\ref{lem:preliminaries-shoelace}, every action term contains a
product of two durations, so \(A\) scales by \(\beta^2\). The integrand in
\(I_K\) is unchanged and the time intervals are all multiplied by \(\beta\), so
\(I_K\) scales by \(\beta\).
\end{proof}

\begin{lemma}[Compactness of simple pure-velocity dual loops]\label{lem:generalized-reeb-compactness}
Let \(z_n\) be simple pure-velocity dual loops whose periods satisfy
\(T_n\to T>0\). After translating the curves and passing to a subsequence,
\(z_n\) converges to a simple pure-velocity dual loop \(z_*\) with period
\(T\), and
\[
  A(z_*)=\lim_n A(z_n),
  \qquad
  I_K(z_*)=T.
\]
\end{lemma}

\begin{proof}
Choose the initial point of each periodic curve at a breakpoint, and insert
zero dwell times for unused facet directions. Each curve is then determined up
to translation by a permutation of the finitely many facet directions and a
nonnegative dwell-time vector of total mass \(T_n\). Pass first to a
subsequence with a fixed permutation. Since \(T_n\to T\), the dwell-time
vectors lie in a common compact set and have a convergent subsequence whose
entries sum to \(T\). Translating the starting points gives uniform convergence
after identifying each domain with \([0,1]\) by linear time rescaling. Closure
passes to the limit because \(\sum_i\tau_iR_{\sigma(i)}=0\) for every curve.
Some limiting dwell times may be zero; deleting those zero-duration entries
recovers the active word of the limiting simple curve. The action is continuous in the dwell-time data by
Lemma~\ref{lem:preliminaries-shoelace}. Since the remaining positive-duration
segments still have pure Reeb velocities, the same computation as in
Lemma~\ref{lem:generalized-reeb-splitting} gives \(I_K(z_*)=T\).
\end{proof}

\begin{proof}[Proof of Theorem~\ref{thm:generalized-reeb-simple-minimizer}]
Choose a capacity-realizing generalized Reeb orbit \(\gamma\). By
Theorem~\ref{thm:preliminaries-clarke-dual-action}, \(\gamma\) is also a dual
minimizer in the thesis convention, with
\[
  I_K(\gamma)=A(\gamma)=T=c_{\mathrm{EHZ}}(K).
\]
Moreover \(\dot\gamma(t)\in\operatorname{conv}\{R_1,\ldots,R_F\}\) for almost
every \(t\).

Choose, by Lemma~\ref{lem:generalized-reeb-pw-affine-approx}, a sequence of
closed piecewise affine curves \(z_{1,n}\to\gamma\) in \(W^{1,2}\), with
velocities in the convex hull of the \(R_i\). Apply
Lemma~\ref{lem:generalized-reeb-splitting} to obtain closed pure-velocity
curves \(z_{2,n}\), and then Lemma~\ref{lem:generalized-reeb-merging} to obtain
closed simple pure-velocity curves \(z_{3,n}\), such that
\[
  A(z_{3,n})\geq A(z_{2,n})\geq A(z_{1,n}),
  \qquad
  I_K(z_{3,n})=T.
\]
Since \(A(z_{1,n})\to A(\gamma)=T\), discard finitely many terms so that
\(A(z_{3,n})>0\). Rescale all durations of \(z_{3,n}\) by
\[
  \beta_n=\frac{T}{A(z_{3,n})},
\]
and denote the resulting curve by \(z_{4,n}\). It satisfies
\[
  A(z_{4,n})=\frac{T^2}{A(z_{3,n})},
  \qquad
  I_K(z_{4,n})=\frac{T^2}{A(z_{3,n})}.
\]
Thus \(z_{4,n}\) is feasible for the dual problem, with period
\(T^2/A(z_{3,n})\). Since \(\gamma\) is a dual minimizer,
\(I_K(z_{4,n})\geq T\), and therefore \(A(z_{3,n})\leq T\). Together with
\(A(z_{3,n})\geq A(z_{1,n})\to T\), this gives
\[
  A(z_{3,n})\to T,
  \qquad \beta_n\to1,
  \qquad \operatorname{period}(z_{4,n})\to T.
\]

By Lemma~\ref{lem:generalized-reeb-compactness}, a subsequence of the \(z_{4,n}\)
converges up to translation to a simple pure-velocity dual loop \(z_*\)
with \(A(z_*)=T\) and \(I_K(z_*)=T\). This curve is a dual minimizer. Applying
Theorem~\ref{thm:preliminaries-clarke-dual-action} again gives a
contact-normalized generalized Reeb orbit on \(\partial K\) in the same
translate/rescale family. In the reconstruction part of that theorem, the
corresponding characteristic has the form
\[
  t\longmapsto c\,z_*(t/c^2)+b/c
\]
for some \(c>0\), with
\[
  c^2=\nu=\frac{I_K(z_*)}{T}=1
\]
by \eqref{eq:preliminaries-dual-multiplier-value}. Thus the reconstruction
only translates \(z_*\). Its velocities
remain the pure facet directions \(R_i\), each used in one connected interval,
so the resulting orbit is simple and has minimum action.  The translation also
recovers the prescribed facet labels.  Indeed, on an interval with
\(\dot\gamma=R_i\), the reconstructed inclusion gives
\[
  2a_i=-J_0R_i
  \in 2\operatorname{conv}\{a_j:\gamma(t)\in F_j\}.
\]
Because the irredundant row \(a_i\) is an extreme point of \(K^\circ\), this
forces \(F_i\) to be active along that interval.  The reconstructed curve lies
in \(K\), so its starting point \(\gamma(0)\) is a valid base point in
Lemma~\ref{lem:generalized-reeb-base-point-recovery}; closure supplies the
remaining displacement equation.
\end{proof}

The theorem supplies one simple capacity minimizer, which is enough to
compute the capacity by a finite optimization. It does not require every
minimizer to be simple. The next section expresses the action in terms of
the word and normalized dwell times, obtaining Haim--Kislev's quadratic
program.
